\documentclass[10pt,a4paper]{amsart}
\usepackage[margin=19mm]{geometry}
\usepackage{amsmath,amssymb,mathtools}
\usepackage[scr=rsfs,cal=euler]{mathalfa}
\usepackage{xfrac}
\usepackage[unicode,hidelinks,bookmarks=false]{hyperref}
\newcommand{\ie}{i.e.\ }
\newcommand{\cf}{cf.\ }
\newcommand{\resp}{respectively\ }
\newcommand{\Subsection}{\subsection}
\providecommand{\nobreakdash}{\nobreak\hbox{-}}
\setlength{\emergencystretch}{2em}
\multlinegap0pt
\usepackage{bm}
\usepackage{bbm}
\usepackage{dsfont}
\usepackage{xspace}
\usepackage{cancel}
\usepackage{mathtools}
\usepackage{amsthm}
\makeatletter
\@ifundefined{lemma}{\newtheorem{lemma}{Lemma}}{}
\@ifundefined{proposition}{\newtheorem{proposition}{Proposition}}{}
\makeatother
\newcommand{\Ker}{\rm{Ker}\>}
\newcommand{\id}{\mathrm{id}}
\newcommand{\im}{\rm{Im}\>}
\title[On the $I(1)$-invariants: Non-abelian Hecke algebra case]{On the $I(1)$-invariants: Non-abelian Hecke algebra case}
\author{Anand Chitrao, Arindam Jana, Asfak Soneji}
\date{Source: arXiv:2604.20218v1 (CC BY 4.0)}
\begin{document}
\maketitle

\noindent\emph{Source and licence.} On the $I(1)$-invariants: Non-abelian Hecke algebra case. Version arXiv:2604.20218v1 (CC BY 4.0), distributed under the Creative Commons Attribution 4.0 International licence, as recorded in the arXiv licence metadata for this exact version and shown on the abstract page https://arxiv.org/abs/2604.20218v1. Full rights notice: licenses/arxiv-2604.20218-CC-BY-4.0.txt.\par\smallskip

\noindent\textit{Editorial scope.} Counted unit: the Proposition whose source label is \texttt{sn1 are invariant} in the article (source0.tex lines 731--736, source label \texttt{sn1 are invariant}), delivered together with the article's own complete original proof of it (source0.tex lines 737--862, one \texttt{proof} environment of 126 lines). No mathematics is written, repaired or re-derived: statement and proof are the article's own text line for line. Required context retained uncounted: snk like terms (lines 592--602); Arindam's identity stolen (lines 707--718). Every definition, display and label of those lines is retained unchanged. Omitted from this package and not redistributed: 1--591, 603--706, 719--730, 863--1759. Nothing inside a retained excerpt is omitted or altered: every retained body is delivered exactly as the article's lines are written, so no omission receipt of any kind accompanies this package. Citations: the retained excerpts carry no citation command at all, so no bibliography entry of the article is transcribed and no reference is redistributed. Reference adaptation: none was needed and none was made; every reference inside the delivered unit resolves inside it, so no unresolved reference or citation is produced. Printed slips kept literal: no printed word, symbol, hypothesis or number of the counted statement or of its proof was altered. Layout and class: the article's own class is \texttt{amsart} with packages including amscd, amssymb, latexsym, amsmath, comment, bbm, amsfonts, tikz-cd, extarrows, mathtools, hyperref, biblatex; the wrapper is the lane's pinned \texttt{amsart} editorial wrapper at 10pt on A4, which restates the theorem-like environments the retained text needs and adds the article's own macro definitions that the retained text needs (\texttt{\textbackslash Ker}, \texttt{\textbackslash id}, \texttt{\textbackslash im}), with extra wrapper packages (bm, bbm, dsfont, xspace, cancel, mathtools). The delivered numbering is the unit's own. Assets: no retained span contains a figure environment, a graphics-inclusion command, a PDF-page inclusion, a table or a TikZ picture; no image file is copied into this package, the package declares no asset dependency and no image file is redistributed with it. Licence: version arXiv:2604.20218v1 of this article is distributed under the Creative Commons Attribution 4.0 International licence, as recorded in the arXiv licence metadata for this exact version and shown on the abstract page https://arxiv.org/abs/2604.20218v1. A full rights notice is delivered at licenses/arxiv-2604.20218-CC-BY-4.0.txt. Characters: every retained excerpt is pure ASCII, so the delivered engine, XeLaTeX with its default Latin Modern text font, reports no missing character.
\begin{lemma}\label{snk like terms}
    Let $n \geq 2$ and $0\leq k \leq q-1$. Then, modulo $({\im} T_{1, 2}, {\Ker} T_{1, 2})$ we have
    \[
        \sum_{\mu \in I_n}c_{[\mu]_{n - 2}}\mu_{n - 2}^{k}\left[g^0_{n, \mu}, 1\right] =
        \begin{cases}
            0 & \text{ if } k \neq q - 1, \\
            \displaystyle\sum_{[\mu]_{n - 2} \in I_{n - 2}}\!\!\!\!\!c_{[\mu]_{n - 2}}[g^0_{n - 2, [\mu]_{n - 2}}, 1] & \text{ if } k = q - 1,
        \end{cases}
    \]
where $c_{[\mu]_{n - 2}} := c(\mu_0, \ldots, \mu_{n - 3})$.
\end{lemma}

    \begin{eqnarray}\label{Arindam's identity stolen}
        \begin{pmatrix}
            1 + \varpi a & b \\ \varpi c & 1 + \varpi d
        \end{pmatrix}
        \begin{pmatrix}
            \varpi & [\lambda] \\ 0 & 1
        \end{pmatrix}
        =
        \begin{pmatrix}
            \varpi & [\lambda + b_0] \\ 0 & 1
        \end{pmatrix}i,
    \end{eqnarray}

\begin{proposition}\label{sn1 are invariant}
Let $e$ and $f$ denote the ramification index and residue degree of the $p$-adic number field $F$, respectively. Assume $ef > 1$ and that $q > 3$. Then, for all $g \in I(1)$, we have 
   $$g \cdot s_{2}^{p^l}-s_{2}^{p^l} \in ({\im}T_{1,2},{\Ker}T_{1,2}),$$
    for all $0 \leq l \leq f-1$. 
    
\end{proposition}

\begin{proof}
  We check the invariance under the following matrices:
$$\begin{pmatrix}
      1+\varpi a  & 0 \\
      0 & 1 
  \end{pmatrix}, \begin{pmatrix}
      1  & 0 \\
      \varpi c & 1
  \end{pmatrix}, \text{ and } \begin{pmatrix}
      1  & b \\
      0 & 1 
  \end{pmatrix},$$
  where $a,b,c \in  \mathcal{O}_F$.
 We begin by computing,
  \begin{eqnarray*}
 \begin{pmatrix}
      1+\varpi a  & 0 \\
      0 & 1 
  \end{pmatrix} \cdot s_{2}^{p^l}  & = & \sum_{\mu \in I_2}\mu_{1}^{p^l}\left[\begin{pmatrix}
      1+\varpi a  & 0 \\
      0 & 1 
  \end{pmatrix} \begin{pmatrix} \varpi^{2} & \mu  \\ 0 & 1\end{pmatrix}, 1\right] \\
  & = &  \sum_{\mu \in I_2}\mu_{1}^{p^l}\left[\begin{pmatrix}
      1+\varpi a  & 0 \\
      0 & 1 
  \end{pmatrix} \begin{pmatrix} \varpi & [\mu_{0}]  \\ 0 & 1\end{pmatrix}\begin{pmatrix} \varpi & [\mu_{1}]  \\ 0 & 1\end{pmatrix}, 1\right]\\
  & = & \sum_{\mu \in I_2}\mu_{1}^{p^l}\left[ \begin{pmatrix} \varpi & [\mu_{0}]  \\ 0 & 1\end{pmatrix}\begin{pmatrix}
      1+\varpi a  & a[\mu_{0}] \\
      0 & 1 
  \end{pmatrix}\begin{pmatrix} \varpi & [\mu_{1}]  \\ 0 & 1\end{pmatrix}, 1\right],
\end{eqnarray*}
which, using \eqref{Arindam's identity stolen}:
$$\begin{pmatrix}
      1+\varpi a  & a[\mu_{0}] \\
      0 & 1 
  \end{pmatrix}\begin{pmatrix} \varpi & [\mu_{1}]  \\ 0 & 1\end{pmatrix}= \begin{pmatrix}
      \varpi & [\mu_{1}+a_{0}\mu_{0}]\\
      0 & 1 
  \end{pmatrix}h,$$
  for some $h \in I(1)$ and making the substitution $\mu_{1}-a_{0}\mu_{0}$, equals %\\
  %$\sum_{\mu \in I_2}(\mu_{1}-a_{0}\mu_{0})^{p^l}\left[ \begin{pmatrix} \varpi & [\mu_{0}]  \\ 0 & 1\end{pmatrix}\begin{pmatrix}
     % \varpi & [\mu_{1}]\\
      %0 & 1 
  %\end{pmatrix}h, 1\right]$
  \begin{eqnarray*}
        %& = & 
        s_{2}^{p^l}+\sum_{\mu \in I_2}(-a_{0}\mu_{0})^{p^l}\left[\begin{pmatrix} \varpi^{2} & \mu  \\ 0 & 1\end{pmatrix}, 1\right]. 
  \end{eqnarray*}
  Applying Lemma \ref{snk like terms}, we get
  $\begin{pmatrix}
      1+\varpi a  & 0 \\
      0 & 1 
  \end{pmatrix} \cdot s_{2}^{p^l} - s_{2}^{p^l} \in ({\im} T_{1,2} , {\Ker} T_{1,2})$ if $q > 2$. If $q = 2$, then this difference is equal to $[\id, 1]$, thereby showing that $s_2^{p^l}$ is not $I(1)$-invariant modulo $({\im} T_{1, 2}, {\Ker} T_{1, 2})$.\\
   We now verify invariance under the lower triangular unipotent elements in $I(1)$:
\begin{eqnarray*}
 \begin{pmatrix}
      1 & 0 \\
      \varpi c & 1 
  \end{pmatrix} \cdot s_{2}^{p^l} & = & \sum_{\mu \in I_2}\mu_{1}^{p^l}\left[\begin{pmatrix}
      1 & 0 \\
      \varpi c & 1 
  \end{pmatrix} \begin{pmatrix} \varpi^{2} & [\mu_{0}]+[\mu_1]\varpi  \\ 0 & 1\end{pmatrix}, 1\right]\\
  &=& \sum_{\mu \in I_2}\mu_{1}^{p^l}\left[\begin{pmatrix}
      1 & 0 \\
      \varpi c & 1 
  \end{pmatrix} \begin{pmatrix} \varpi & [\mu_{0}]  \\ 0 & 1\end{pmatrix}\begin{pmatrix} \varpi & [\mu_1]  \\ 0 & 1\end{pmatrix}, 1\right]\\
  &=& \sum_{\mu \in I_2}\mu_{1}^{p^l}\left[ \begin{pmatrix} \varpi & [\mu_{0}]  \\ 0 & 1\end{pmatrix} \begin{pmatrix}
      1-\varpi c[\mu_{0}] & -[\mu_{0}^2]c \\
      \varpi^2 c & 1+\varpi c[\mu_{0}] 
  \end{pmatrix}\begin{pmatrix} \varpi & [\mu_1]  \\ 0 & 1\end{pmatrix}, 1\right]\\
  &=& \sum_{\mu \in I_2}\mu_{1}^{p^l}\left[ \begin{pmatrix} \varpi & [\mu_{0}]  \\ 0 & 1\end{pmatrix} \begin{pmatrix}
      \varpi & [\mu_{1}-\mu_{0}^2c_{0}] \\
      0 & 1 
  \end{pmatrix}i, 1\right],
\end{eqnarray*}
for some $i \in I(1)$, by \eqref{Arindam's identity stolen}. Making the change of variable $\mu_{1} \mapsto \mu_{1}+c_{0}\mu_{0}^2$, we obtain
\begin{eqnarray*}
  \begin{pmatrix}
      1 & 0 \\
      \varpi c & 1 
  \end{pmatrix} \cdot s_{2}^{p^l}&=&\sum_{\mu \in I_2}(\mu_{1}+c_{0}\mu_{0}^2)^{p^l}\left[\begin{pmatrix} \varpi^{2} & \mu  \\ 0 & 1\end{pmatrix}, 1 \right]\\ 
  &=& s_{2}^{p^l}+\sum_{\mu \in I_2}(c_{0}\mu_{0}^2)^{p^l}\left[\begin{pmatrix} \varpi^{2} & \mu  \\ 0 & 1\end{pmatrix}, 1 \right].
\end{eqnarray*}
Applying Lemma \ref{snk like terms} once more, we conclude (for $q > 3$) that
\[
\begin{pmatrix}
      1 & 0 \\
      \varpi c & 1 
  \end{pmatrix} \cdot s_{2}^{p^l}-s_{2}^{p^l} \in ({\im} T_{1,2} , {\Ker} T_{1,2}) .
\]
Now, we compute
\begin{eqnarray*}
 \begin{pmatrix}
      1 & b \\
      0 & 1 
  \end{pmatrix} \cdot s_{2}^{p^l} & = & \sum_{\mu \in I_2}\mu_{1}^{p^l}\left[\begin{pmatrix}
      1 & b \\
      0 & 1 
  \end{pmatrix} \begin{pmatrix} \varpi^{2} & [\mu_{0}]+[\mu_1]\varpi  \\ 0 & 1\end{pmatrix}, 1\right]\\
  &=& \sum_{\mu \in I_2}\mu_{1}^{p^l}\left[\begin{pmatrix} \varpi^{2} & [\mu_{0}]+[\mu_1]\varpi+b  \\ 0 & 1\end{pmatrix}, 1\right]\\
  &=& \sum_{\mu \in I_2}\mu_{1}^{p^l}\left[\begin{pmatrix} \varpi & [\mu_{0}+b_0]  \\ 0 & 1\end{pmatrix}\begin{pmatrix} \varpi & [\mu_1+b_1+Z]  \\ 0 & 1\end{pmatrix}i, 1\right],
\end{eqnarray*}
for some $i \in I(1)$, where the last equality follows from \eqref{Arindam's identity stolen}. Here
\[
Z=\begin{cases}
      0, &  \text{if } e>1, \\
      \dfrac{\mu_0^{q}+b_0^{q}-(\mu_{0}+b_0)^q}{p}, & \text{if } e=1.
  \end{cases}
\] \\
Next, we perform the change of variables
$$\mu_{0} \rightarrow \mu_{0}-b_0 \quad  \text{and} \quad \mu_{1} \rightarrow \mu_{1}-b_1-Z,$$
which yields
\begin{eqnarray*}
 \sum_{\mu \in I_2}(\mu_{1}-b_1-Z)^{p^l}\left[\begin{pmatrix} \varpi^2 & \mu  \\ 0 & 1\end{pmatrix}, 1\right]
 &= & s_{2}^{p^l}+\sum_{\mu \in I_2}(-b_1-Z)^{p^l}\left[\begin{pmatrix} \varpi^2 & \mu  \\ 0 & 1\end{pmatrix}, 1\right].
\end{eqnarray*}
If $e > 1$, then $Z = 0$ and therefore the largest power of $\mu_0$ appearing in the coefficient is at most $0$. If $e = 1$, then $f > 1$ and therefore the largest power of $\mu_0$ appearing in the coefficient is $(p - 1)p^{f - 1} < q - 1$. In both cases, Lemma~\ref{snk like terms} applies, and we get
%Applying Lemma \ref{snk like terms}, we obtain
\[
\begin{pmatrix}
      1 & b \\
      0 & 1
\end{pmatrix} \cdot s_{2}^{p^l}-s_2^{p^l} \in ({\im} T_{1,2}, {\Ker} T_{1,2}).
\]
This completes the proof of the proposition.
\end{proof}

\end{document}
